\documentclass[11pt]{article}
\usepackage[margin=2.4cm]{geometry}
\usepackage{amsmath,amssymb,booktabs,graphicx,xspace,microtype,longtable}
\usepackage{alltt}
% cmtt has no OMS shape, so a maths-font backslash was being substituted
% inside alltt and did not match the verbatim blocks. This one comes from
% whatever font the block is already set in.
\newcommand{\bs}{\char92\relax}
\usepackage[hyphens]{url}
\usepackage[colorlinks=true,allcolors=blue]{hyperref}
% Long literals are single unbreakable words. \path lets a line break at the
% punctuation a path already contains; a SHA-256 digest contains none, so
% \digest also permits a break after a hexadecimal character. Neither command
% adds or removes a character: only the line breaking changes.
\DeclareUrlCommand\digest{\urlstyle{tt}%
  \def\UrlBreaks{\do\0\do\1\do\2\do\3\do\4\do\5\do\6\do\7\do\8\do\9%
    \do\a\do\b\do\c\do\d\do\e\do\f}}
% Generated by tools/make_macros.py -- do not edit by hand.
% Every numeric result in the manuscript comes from this file.
\newcommand{\ABcases}{10\,000\xspace}
\newcommand{\Acases}{98\,385\xspace}
\newcommand{\Achecks}{984\,213\xspace}
\newcommand{\Acomp}{363\xspace}
\newcommand{\Afailed}{0\xspace}
\newcommand{\Amaxlen}{8\xspace}
\newcommand{\Awords}{9\,840\xspace}
\newcommand{\AwordsFull}{9\,828\xspace}
\newcommand{\AwordsShortOne}{3\xspace}
\newcommand{\AwordsShortTwo}{9\xspace}
\newcommand{\BOneFootprint}{0\xspace}
\newcommand{\BOneInterior}{10\,000\xspace}
\newcommand{\BThreeFootprint}{0\xspace}
\newcommand{\BThreeInterior}{0\xspace}
\newcommand{\BTwoFootprint}{10\,000\xspace}
\newcommand{\BTwoInterior}{0\xspace}
\newcommand{\BZeroFootprint}{10\,000\xspace}
\newcommand{\BZeroInterior}{10\,000\xspace}
\newcommand{\BbaseFive}{9\,455\xspace}
\newcommand{\BbaseLineageOne}{10\,000\xspace}
\newcommand{\BbaseOne}{9\,455\xspace}
\newcommand{\BbaseOnePct}{94.5\xspace}
\newcommand{\BbaseSupportOne}{0\xspace}
\newcommand{\BbaseUnverOne}{7\,222\xspace}
\newcommand{\BctrlDev}{0.0070\xspace}
\newcommand{\BctrlDevPct}{0.70\xspace}
\newcommand{\Bdepths}{5\xspace}
\newcommand{\BemAll}{0\xspace}
\newcommand{\BemLineage}{0\xspace}
\newcommand{\BemOpAll}{0\xspace}
\newcommand{\BemUnver}{0\xspace}
\newcommand{\Bintervals}{64\xspace}
\newcommand{\BopCount}{9\xspace}
\newcommand{\BopMax}{100.0\xspace}
\newcommand{\BopMin}{94.3\xspace}
\newcommand{\Btimelines}{10\,000\xspace}
\newcommand{\CaptSupport}{0.90\xspace}
\newcommand{\Cbase}{600\xspace}
\newcommand{\CbaseLineage}{600\xspace}
\newcommand{\Cclips}{600\xspace}
\newcommand{\CcorpusFs}{16\,000\xspace}
\newcommand{\CcorpusMismatch}{0\xspace}
\newcommand{\Cem}{0\xspace}
\newcommand{\CemLineage}{0\xspace}
\newcommand{\Cexact}{600\xspace}
\newcommand{\Cgen}{600\xspace}
\newcommand{\Clabel}{io.github.1d-coda.emaudio.evidence\xspace}
\newcommand{\CmeanDurS}{2.85\xspace}
\newcommand{\Cnamespace}{io.github.1d-coda\xspace}
\newcommand{\Cschema}{https://github.com/1D-Coda/em-audio/blob/main/docs/em-audio-schema-1.0.md\xspace}
\newcommand{\CtwoNeuralBase}{50\xspace}
\newcommand{\CtwoNeuralEm}{0\xspace}
\newcommand{\CtwoNeuralExact}{50\xspace}
\newcommand{\CtwoNeuralN}{50\xspace}
\newcommand{\CtwoNoiseEm}{0\xspace}
\newcommand{\CtwoNoiseExact}{50\xspace}
\newcommand{\CtwoNoiseN}{50\xspace}
\newcommand{\Dbase}{3\,600\xspace}
\newcommand{\DbaseLineage}{4\,200\xspace}
\newcommand{\DbasePromo}{3\,600\xspace}
\newcommand{\DdeclRejected}{1\,223\xspace}
\newcommand{\Ddeterm}{0\xspace}
\newcommand{\DdetermN}{100\xspace}
\newcommand{\Deither}{4\,200\xspace}
\newcommand{\DeitherPct}{87.5\xspace}
\newcommand{\Dem}{0\xspace}
\newcommand{\DemInterval}{0\xspace}
\newcommand{\DemLineage}{0\xspace}
\newcommand{\DemPromo}{0\xspace}
\newcommand{\DfallbackOutputs}{888\xspace}
\newcommand{\DfallbackSamples}{187\,902\xspace}
\newcommand{\DguardOK}{all\xspace}
\newcommand{\Dinterval}{4\,200\xspace}
\newcommand{\DmaxDev}{681\xspace}
\newcommand{\DmeanEmIv}{4.19\xspace}
\newcommand{\DmpThreeDev}{31\xspace}
\newcommand{\DoverlayInterval}{600\xspace}
\newcommand{\DresampleDev}{0\xspace}
\newcommand{\Druns}{4\,800\xspace}
\newcommand{\DselInterval}{0\xspace}
\newcommand{\DshorterOut}{335\xspace}
\newcommand{\DsilenceDev}{681\xspace}
\newcommand{\DstretchDev}{471\xspace}
\newcommand{\DstrictLineage}{0\xspace}
\newcommand{\DstrictPromo}{0\xspace}
\newcommand{\DstrictVerified}{4\,800\xspace}
\newcommand{\DtotalRuns}{4\,800\xspace}
\newcommand{\Dtransforms}{8\xspace}
\newcommand{\DzeroBase}{2\xspace}
\newcommand{\Eclips}{60\xspace}
\newcommand{\Econds}{5\xspace}
\newcommand{\Ederived}{60\xspace}
\newcommand{\Ereenc}{60\xspace}
\newcommand{\Estripped}{60\xspace}
\newcommand{\Etampered}{60\xspace}
\newcommand{\Evalid}{60\xspace}
\newcommand{\EvalidDerived}{60\xspace}
\newcommand{\Eviol}{0\xspace}
\newcommand{\Fassertion}{120\xspace}
\newcommand{\FassertionFlac}{1\,334\xspace}
\newcommand{\FassertionMpThree}{2\,114\xspace}
\newcommand{\FassertionWav}{1\,333\xspace}
\newcommand{\Fclips}{40\xspace}
\newcommand{\Fcontainerlist}{\texttt{wav}, \texttt{flac} and \texttt{mp3}\xspace}
\newcommand{\Fcontainers}{3\xspace}
\newcommand{\Fderived}{120\xspace}
\newcommand{\Ffilechanged}{120\xspace}
\newcommand{\Fmismatch}{0\xspace}
\newcommand{\FmismatchPolicy}{0\xspace}
\newcommand{\Fparent}{120\xspace}
\newcommand{\Ftotal}{120\xspace}
\newcommand{\Ftrusted}{120\xspace}
\newcommand{\Fvalfail}{0\xspace}
\newcommand{\Gassertion}{1\,333\xspace}
\newcommand{\Gbase}{0.41\xspace}
\newcommand{\Gclips}{30\xspace}
\newcommand{\Gem}{1.35\xspace}
\newcommand{\GenSupport}{0.10\xspace}
\newcommand{\GffmpegPct}{0.152\xspace}
\newcommand{\GfirstIterPct}{0.09\xspace}
\newcommand{\Gmanifest}{14\,932\xspace}
\newcommand{\GmaxIntervals}{256\xspace}
\newcommand{\GperInterval}{333\xspace}
\newcommand{\GperIntervalUs}{21.0\xspace}
\newcommand{\Gratio}{3.34\xspace}
\newcommand{\Greps}{30\xspace}
\newcommand{\Gsign}{16.9\xspace}
\newcommand{\GsignQa}{16.2\xspace}
\newcommand{\GsignQb}{17.2\xspace}
\newcommand{\Gval}{14.0\xspace}
\newcommand{\GvalQa}{13.6\xspace}
\newcommand{\GvalQb}{14.4\xspace}
\newcommand{\Hcases}{10\,860\xspace}
\newcommand{\Hdelta}{0\xspace}
\newcommand{\Hdis}{0\xspace}
\newcommand{\IRabs}{1.62\xspace}
\newcommand{\IRabsCv}{13.2\xspace}
\newcommand{\IRabsGap}{20\xspace}
\newcommand{\IRabsRef}{1.35\xspace}
\newcommand{\IRachecks}{885\,828\xspace}
\newcommand{\IRafailed}{0\xspace}
\newcommand{\IRarch}{x86\_64\xspace}
\newcommand{\IRcTwopatool}{0.27.15\xspace}
\newcommand{\IRclean}{9\xspace}
\newcommand{\IRcores}{12\xspace}
\newcommand{\IRcpu}{13th Gen Intel Core i7-1355U\xspace}
\newcommand{\IRdate}{2026-08-27\xspace}
\newcommand{\IRdet}{216\xspace}
\newcommand{\IRdetBuild}{213\xspace}
\newcommand{\IRdetSuite}{3\xspace}
\newcommand{\IRdiffFiles}{4\xspace}
\newcommand{\IRenv}{28\xspace}
\newcommand{\IRespeak}{1.52.0\xspace}
\newcommand{\IRffmpeg}{8.0.1-3ubuntu2\xspace}
\newcommand{\IRfiles}{13\xspace}
\newcommand{\IRmem}{8\xspace}
\newcommand{\IRmpContext}{transient\xspace}
\newcommand{\IRmpDeclared}{2\,304\xspace}
\newcommand{\IRmpExcess}{2\,013\xspace}
\newcommand{\IRmpOutside}{110\xspace}
\newcommand{\IRmpReach}{4\,317\xspace}
\newcommand{\IRmpReachRef}{1\,555\xspace}
\newcommand{\IRnode}{v22.22.1\xspace}
\newcommand{\IRopsAgree}{3\xspace}
\newcommand{\IRopsNonzero}{4\xspace}
\newcommand{\IRopsZero}{3\xspace}
\newcommand{\IRos}{Linux 6.6.87.2\xspace}
\newcommand{\IRpython}{3.14.4\xspace}
\newcommand{\IRratio}{3.54\xspace}
\newcommand{\IRratioCv}{1.6\xspace}
\newcommand{\IRratioGap}{5.9\xspace}
\newcommand{\IRratioRef}{3.34\xspace}
\newcommand{\IRsilDev}{12\,505\xspace}
\newcommand{\IRsilDevRef}{681\xspace}
\newcommand{\IRsilGuard}{4\,096\xspace}
\newcommand{\IRsilLen}{16\,384\xspace}
\newcommand{\IRsilLenRef}{13\,312\xspace}
\newcommand{\IRsilMedian}{9\,105.5\xspace}
\newcommand{\IRtag}{v1.0.0\xspace}
\newcommand{\IchainClips}{200\xspace}
\newcommand{\IchainDeep}{85.74\xspace}
\newcommand{\IchainDeepMax}{99.99\xspace}
\newcommand{\IchainDepth}{5\xspace}
\newcommand{\IclipsPerTf}{600\xspace}
\newcommand{\IemMaxMedianPct}{31.34\xspace}
\newcommand{\IlongShortPct}{25.13\xspace}
\newcommand{\IlongThirtyPct}{2.40\xspace}
\newcommand{\IlongThreeHundredPct}{0.240\xspace}
\newcommand{\ImaxAnyPct}{58.33\xspace}
\newcommand{\ImpThreeMedianPct}{20.24\xspace}
\newcommand{\IoverlayMedianPct}{0.00\xspace}
\newcommand{\IresampleMedianPct}{0.86\xspace}
\newcommand{\IsilenceMedianPct}{31.34\xspace}
\newcommand{\IstretchMedianPct}{22.64\xspace}
\newcommand{\IstrictFlacPct}{0.00\xspace}
\newcommand{\IstrictMpPct}{100.00\xspace}
\newcommand{\IstrictSelPct}{100.00\xspace}
\newcommand{\IstrictStPct}{99.99\xspace}
\newcommand{\ItfCount}{8\xspace}
\newcommand{\IwholeMinMedianPct}{99.99\xspace}
\newcommand{\IwholeOverlayPct}{96.99\xspace}
\newcommand{\IwholePromoted}{0\xspace}
\newcommand{\IzeroTf}{4\xspace}
\newcommand{\JcompTrusted}{30\xspace}
\newcommand{\Jcomponents}{30\xspace}
\newcommand{\JderivedMixed}{30\xspace}
\newcommand{\JderivedTrusted}{30\xspace}
\newcommand{\Jessence}{0\xspace}
\newcommand{\Jfixtures}{30\xspace}
\newcommand{\Jmixed}{30\xspace}
\newcommand{\Jplaced}{30\xspace}
\newcommand{\Jroundtrip}{30\xspace}
\newcommand{\JsourcesTrusted}{60\xspace}
\newcommand{\Kaffected}{52\,005\xspace}
\newcommand{\KalignResample}{1\xspace}
\newcommand{\KalignSelection}{205\xspace}
\newcommand{\KalignStretch}{401\xspace}
\newcommand{\KconfTotal}{0\xspace}
\newcommand{\Kcontexts}{4\xspace}
\newcommand{\KfallbackTotal}{15\,390\xspace}
\newcommand{\KimpulseN}{16\,384\xspace}
\newcommand{\KmpThreeFp}{2\,304\xspace}
\newcommand{\KmpThreeMargin}{749\xspace}
\newcommand{\KmpThreeReach}{1\,555\xspace}
\newcommand{\KmpThreeSpread}{1\,823\xspace}
\newcommand{\Kops}{7\xspace}
\newcommand{\Koutside}{0\xspace}
\newcommand{\KoutsideTotal}{0\xspace}
\newcommand{\Kprobes}{280\xspace}
\newcommand{\KresampleFp}{97\xspace}
\newcommand{\KresampleGuard}{64\xspace}
\newcommand{\KresampleKernel}{33\xspace}
\newcommand{\KresampleMargin}{65\xspace}
\newcommand{\KresampleReach}{32\xspace}
\newcommand{\KresampleSpread}{33\xspace}
\newcommand{\KsilenceFp}{4\,096\xspace}
\newcommand{\KsilenceMargin}{819\xspace}
\newcommand{\KsilenceReach}{3\,277\xspace}
\newcommand{\KstretchAboveOne}{7\xspace}
\newcommand{\KstretchFp}{2\,577\xspace}
\newcommand{\KstretchMargin}{671\xspace}
\newcommand{\KstretchReach}{1\,906\xspace}
\newcommand{\KstretchSpread}{2\,023\xspace}
\newcommand{\KstrictProfileTotal}{0\xspace}
\newcommand{\KstrictTotal}{6\xspace}
\newcommand{\KtwoAmpMax}{12\,000\xspace}
\newcommand{\KtwoAmpMin}{256\xspace}
\newcommand{\KtwoAmps}{4\xspace}
\newcommand{\KtwoContentExceeding}{2\xspace}
\newcommand{\KtwoContexts}{6\xspace}
\newcommand{\KtwoDemoMpDeclared}{captured\xspace}
\newcommand{\KtwoDemoMpOut}{7\,355\xspace}
\newcommand{\KtwoDemoMpSrc}{900\xspace}
\newcommand{\KtwoDemoMpStrict}{mixed\xspace}
\newcommand{\KtwoDemoStDeclared}{captured\xspace}
\newcommand{\KtwoDemoStOut}{3\,162\xspace}
\newcommand{\KtwoDemoStSrc}{900\xspace}
\newcommand{\KtwoDemoStStrict}{mixed\xspace}
\newcommand{\KtwoHoldoutContexts}{2\xspace}
\newcommand{\KtwoMpExceed}{4\xspace}
\newcommand{\KtwoMpReach}{11\,292\xspace}
\newcommand{\KtwoNormCtrlExceed}{84\xspace}
\newcommand{\KtwoNormCtrlReach}{15\,600\xspace}
\newcommand{\KtwoPositions}{8\xspace}
\newcommand{\KtwoProbes}{1\,152\xspace}
\newcommand{\KtwoProbesPerOp}{192\xspace}
\newcommand{\KtwoResReach}{32\xspace}
\newcommand{\KtwoResampleReach}{32\xspace}
\newcommand{\KtwoSelReach}{3\,277\xspace}
\newcommand{\KtwoStExceed}{14\xspace}
\newcommand{\KtwoStMinAmp}{256\xspace}
\newcommand{\KtwoStReach}{15\,482\xspace}
\newcommand{\KtwoStrictProbes}{0\xspace}
\newcommand{\KunmappedTotal}{6\xspace}
\newcommand{\LscopeCases}{441\xspace}
\newcommand{\LscopeFailBad}{80\xspace}
\newcommand{\LscopeScopes}{5\xspace}
\newcommand{\LscopeSkipped}{184\xspace}
\newcommand{\LscopeViolations}{0\xspace}
\newcommand{\MpThreeDeclared}{2\,304\xspace}
\newcommand{\MpThreeGuard}{576\xspace}
\newcommand{\MpThreeKernel}{1\,728\xspace}
\newcommand{\MpThreeWindow}{1\,152\xspace}
\newcommand{\MstabAbsPct}{2.14\xspace}
\newcommand{\MstabFactor}{1.1\xspace}
\newcommand{\MstabRatioPct}{2.05\xspace}
\newcommand{\MstabRepeats}{9\xspace}
\newcommand{\NoteComposition}{S5\xspace}
\newcommand{\NoteHoldout}{S9\xspace}
\newcommand{\NoteIndependent}{S8\xspace}
\newcommand{\NoteNovelty}{S2\xspace}
\newcommand{\NoteRawRows}{S10\xspace}
\newcommand{\NoteRepro}{S7\xspace}
\newcommand{\NoteRobustness}{S6\xspace}
\newcommand{\NoteSchema}{S4\xspace}
\newcommand{\NoteThreat}{S3\xspace}
\newcommand{\NoteTransfer}{S1\xspace}
\newcommand{\PbaseFlagged}{3\,600\xspace}
\newcommand{\PbaseLocal}{874\xspace}
\newcommand{\PbasePromo}{3\,600\xspace}
\newcommand{\PbytesDecl}{2\,584\xspace}
\newcommand{\PbytesPlain}{2\,093\xspace}
\newcommand{\PgapCases}{4\,200\xspace}
\newcommand{\PgapFlagged}{4\,200\xspace}
\newcommand{\PhonestDecl}{4\,800\xspace}
\newcommand{\PhonestStrict}{4\,800\xspace}
\newcommand{\PjsCases}{2\,400\xspace}
\newcommand{\PjsDis}{0\xspace}
\newcommand{\PlinCases}{3\,000\xspace}
\newcommand{\PlinFlagged}{3\,000\xspace}
\newcommand{\PnarrowCases}{2\,400\xspace}
\newcommand{\PnarrowPassed}{2\,400\xspace}
\newcommand{\Pout}{4\,800\xspace}
\newcommand{\PropFootprint}{2\xspace}
\newcommand{\PropUnion}{1\xspace}
\newcommand{\PshortCases}{4\,800\xspace}
\newcommand{\PshortFlagged}{4\,800\xspace}
\newcommand{\PtamperCases}{2\,804\xspace}
\newcommand{\PtamperFlagged}{2\,804\xspace}
\newcommand{\PtwoClips}{20\xspace}
\newcommand{\PtwoHonest}{20\xspace}
\newcommand{\PtwoWrongFlagged}{20\xspace}
\newcommand{\PtwoWrongValid}{20\xspace}
\newcommand{\PverifyMs}{0.16\xspace}
\newcommand{\Qbase}{4\,006\xspace}
\newcommand{\Qclips}{600\xspace}
\newcommand{\QconsFlagged}{4\,006\xspace}
\newcommand{\QconsHonest}{4\,800\xspace}
\newcommand{\QconsLocal}{625\xspace}
\newcommand{\QconsPromo}{4\,006\xspace}
\newcommand{\Qeither}{4\,158\xspace}
\newcommand{\QeitherPct}{86.6\xspace}
\newcommand{\Qem}{0\xspace}
\newcommand{\QemInterval}{0\xspace}
\newcommand{\QemLineage}{0\xspace}
\newcommand{\QfallbackOutputs}{1\,137\xspace}
\newcommand{\QgenPct}{58.6\xspace}
\newcommand{\Qinterval}{4\,158\xspace}
\newcommand{\QmarginExceeded}{1\xspace}
\newcommand{\QmedDur}{3.45\xspace}
\newcommand{\QmedIntervals}{5\xspace}
\newcommand{\QoverlayDev}{5\xspace}
\newcommand{\Qruns}{4\,800\xspace}
\newcommand{\QspoofedDev}{15\,847\xspace}
\newcommand{\Qstep}{26\xspace}
\newcommand{\QstretchDev}{965\xspace}
\newcommand{\Qstrict}{0\xspace}
\newcommand{\QstrictVerified}{4\,800\xspace}
\newcommand{\Rbase}{3\,710\xspace}
\newcommand{\Rclips}{600\xspace}
\newcommand{\Rcommit}{ceaf661e5f7a0e8ae7ff00dd29f7fd343d3a8858\xspace}
\newcommand{\RconsFlagged}{3\,710\xspace}
\newcommand{\RconsPromo}{3\,710\xspace}
\newcommand{\Reither}{4\,315\xspace}
\newcommand{\Rem}{0\xspace}
\newcommand{\RemInterval}{0\xspace}
\newcommand{\RemLineage}{0\xspace}
\newcommand{\RfallbackOutputs}{826\xspace}
\newcommand{\RfilesTotal}{42\,471\xspace}
\newcommand{\RgenPct}{19.0\xspace}
\newcommand{\RjsCases}{2\,400\xspace}
\newcommand{\RjsDis}{0\xspace}
\newcommand{\RmarginExceeded}{0\xspace}
\newcommand{\RmedDur}{3.55\xspace}
\newcommand{\RmedIntervals}{3\xspace}
\newcommand{\RperSpk}{20\xspace}
\newcommand{\Rruns}{4\,800\xspace}
\newcommand{\Rspeakers}{30\xspace}
\newcommand{\Rstrict}{0\xspace}
\newcommand{\RstrictVerified}{4\,800\xspace}
\newcommand{\Rtag}{v1.3.1\xspace}
\newcommand{\SecComposition}{7.9\xspace}
\newcommand{\SecHoldout}{7.8\xspace}
\newcommand{\SecIndependent}{7.14\xspace}
\newcommand{\SecOverhead}{7.13\xspace}
\newcommand{\TabOperators}{2\xspace}
\newcommand{\Tfail}{0\xspace}
\newcommand{\ThmComposition}{1\xspace}
\newcommand{\Tpass}{49\xspace}
\newcommand{\Ttotal}{49\xspace}
\newcommand{\Varch}{arm64\xspace}
\newcommand{\VcTwopatool}{0.27.2\xspace}
\newcommand{\Vcores}{14\xspace}
\newcommand{\Vcpu}{Apple M3 Max\xspace}
\newcommand{\Vespeak}{1.52.0\xspace}
\newcommand{\Vffmpeg}{9.0.1\xspace}
\newcommand{\Vnode}{v26.7.0\xspace}
\newcommand{\Vos}{macOS 26.5.2\xspace}
\newcommand{\Vpython}{3.11.15\xspace}
\newcommand{\Vram}{36\xspace}
\newcommand{\XBMacChecks}{885\,828\xspace}
\newcommand{\XBMacCpu}{Apple M1\xspace}
\newcommand{\XBMacEnvDiffs}{17\xspace}
\newcommand{\XBMacFailed}{0\xspace}
\newcommand{\XBMacFfmpeg}{9.0.1\xspace}
\newcommand{\XBMacOs}{macOS 26.3.1\xspace}
\newcommand{\XBMacOutside}{0\xspace}
\newcommand{\XBMacReach}{1\,555\xspace}
\newcommand{\XBMacVerdict}{every deterministic output\xspace}
\newcommand{\XBRefOutside}{0\xspace}
\newcommand{\XBRefReach}{1\,555\xspace}
\newcommand{\XBRefStretchDev}{471\xspace}
\newcommand{\XBStretchGuard}{2\,048\xspace}
\newcommand{\XBWinChecks}{984\,213\xspace}
\newcommand{\XBWinCpu}{Intel64 Family 6 Model 186 Stepping 3, GenuineIntel\xspace}
\newcommand{\XBWinDetDiffs}{69\xspace}
\newcommand{\XBWinFailed}{0\xspace}
\newcommand{\XBWinFfmpeg}{9.0.1-full\_build-www.gyan.dev\xspace}
\newcommand{\XBWinOs}{Windows 11\xspace}
\newcommand{\XBWinOutside}{0\xspace}
\newcommand{\XBWinReach}{1\,541\xspace}
\newcommand{\XBWinStretchDev}{463\xspace}
\newcommand{\XBWinTag}{v1.0.7\xspace}

\newcommand{\C}{\ensuremath{C}}
\newcommand{\G}{\ensuremath{G}}
\renewcommand{\thesection}{S\arabic{section}}
\renewcommand{\thetable}{S\arabic{table}}

\title{Supplementary material for\\
\emph{Evidence-Monotone Audio Representation: Preventing Provenance Promotion
Across Derived Audio Transformations}}
\author{Alejandro Urias}
\date{}

\begin{document}
\maketitle

\section{Transfer memo: what is prior work and what is new}
\label{s:transfer}

Every element in the first column is prior work and is cited rather than
claimed; every reported number in this paper is new.

\begin{table}[htbp]
\centering
\caption{Relation to the spatial predecessor, cited in the manuscript.}
\label{tab:inherited}
\small
\begin{tabular}{p{0.46\linewidth}p{0.46\linewidth}}
\toprule
Inherited from the spatial work & New in this work \\
\midrule
Evidence record $(P,S,A,L)$; provenance as an atom set with mixed by union &
Temporal required-source sets over half-open sample intervals; the
captured/generated alphabet and the capture-support channel \\
\addlinespace
Same-channel meet by minimum; scope intersection; lineage union &
Requirement (iv): a channel not applicable to \emph{every} required source is
withheld, not computed from the rest \\
\addlinespace
Absorbing unverified state &
Requirement (iv$'$): a $\bot$ source also withholds every numeric channel \\
\addlinespace
Non-promotion in the spatial setting &
Restated without reference to a modality, as an instance of a known
meet-semilattice construction; kernel footprints with declared source-support
bounds, validated against stock FFmpeg, tested by differential influence probes
and challenged on held-out inputs; footprint monotonicity
(Proposition~\PropFootprint{}) \\
\addlinespace
Conformance-test method (exhaustive enumeration, frozen fixtures, separately
implemented oracle) &
Signal transparency on decoded PCM; signed C2PA transport with temporal
regions; the claim-dilution metric; the strict profile and the consumer
verifier \\
\bottomrule
\end{tabular}
\end{table}

\section{Novelty search and closest-work matrix}
\label{s:novelty}

\paragraph{Search date} First pass 2026-08-19; re-run 2026-09-05 before
submission. The re-run is reported below with the first pass, since what a
novelty search establishes depends on when it was last attempted.

\paragraph{Databases and interfaces} Open-literature web search; arXiv; IACR
Cryptology ePrint Archive; ACM Digital Library and LIPIcs listings; the C2PA
specification site; OpenSLR.

\begingroup
% The query strings are set in typewriter, whose interword space has no
% stretch, so justifying them stretches the surrounding roman text.
\raggedright
\paragraph{Query strings}
(1) \texttt{provenance semiring non-amplification derived media operator
monotone provenance lattice};
(2) \texttt{C2PA content credentials audio provenance derived asset ingredient
inheritance research paper 2026};
(3) \texttt{"provenance laundering" OR "metadata washing" media provenance
attack C2PA};
(4) \texttt{audio editing provenance preservation composition proof interval
lineage trim concatenate resample formal};
(5) \texttt{Denning lattice model secure information flow non-interference
applied to media provenance labels 2025 2026};
(6) \texttt{audio deepfake detection survey 2026 speech antispoofing ASVspoof 5
generalization};
(7) \texttt{audio watermarking soft binding C2PA durable content credentials
AudioSeal 2026}.
\par\endgroup

\paragraph{Outcome} The columns \emph{complete-source operator rule} and
\emph{C2PA-integrated interval-level enforcement} were empty across every prior
row. Because provenance-semiring and information-flow work already yields the
non-amplification property abstractly, the theorem is not claimed as new; it is
presented as an instance and the contribution is relocated to the executable
temporal contract, the transport and the measurements. The full matrix is in
\path{paper/novelty_search.md} in the repository.

\paragraph{Outcome of the re-run} The re-run was conducted to falsify the
novelty claim rather than to support it, and targeted the three statements this
work depends on: that a derived segment must not carry provenance stronger than
its weakest required source; that per-operator kernel support is declared and
calibrated rather than assumed; and that the resulting interval evidence is
carried in a signed C2PA manifest over temporal regions. Within the searched
sources, no prior system was found that combines conservative evidence
reduction, explicit temporal dependency declarations and a signed C2PA audio
transport; the first statement on its own is the meet-semilattice reading of
the cited lattice literature and of the spatial predecessor. The searches did return the established algebraic
material the paper takes as given, which is the outcome the claim predicts: the
meet-semilattice construction and the restriction of semiring provenance to
monotone queries are prior art and are cited as such. Two 2026 preprints are the
closest adjacent work, one on anonymous provenance for verifiably edited audio
and one a security analysis of C2PA and its implementations; both were already
cited before the re-run, and neither defines an inheritance rule over media
time. C2PA itself records multiple sources as ingredients but specifies no rule
constraining what a derived asset may claim relative to them; this work
supplies that rule. Nothing found in the re-run changed a claim in this paper
or added a reference; a search that surfaces only works already cited is the
result we report, not evidence that none was sought. A later journal-context
search conducted during target-journal preparation added four FSI:DI papers on
forensic representation and tool-output variability for problem framing; none
altered the novelty claims or the closest-work assessment above.

\section{Threat-model matrix}
\label{s:threat}

\begin{longtable}{p{\dimexpr0.30\linewidth-2\tabcolsep\relax}%
                  p{\dimexpr0.36\linewidth-2\tabcolsep\relax}%
                  p{\dimexpr0.33\linewidth-2\tabcolsep\relax}}
\caption{Threat-model matrix: what each attacker capability meets in the
contract, and what the contract does not prevent.}\label{tab:threat}\\
\toprule
Attacker capability & Contract behaviour & Not prevented \\
\midrule
\endfirsthead
\toprule
Attacker capability & Contract behaviour & Not prevented \\
\midrule
\endhead
Passes a mixed-origin asset through a benign transformation to obtain a
captured-only claim &
Output claim is the meet over the complete required source set, so it stays
mixed &
Nothing, if the attacker also controls the signer \\
\addlinespace
Trims away the boundaries so that only captured material is at the edges &
Interior sources remain in $D_y$; claim stays mixed &
A trim that genuinely removes all generated material, which legitimately yields
a captured-only claim \\
\addlinespace
Relies on a filtered operator so that a generated sample bleeds into a
neighbouring output region without appearing in the nominal range &
The declared kernel footprint enlarges $D_y$; over-approximation is safe by
Proposition~\PropFootprint{} (footprint monotonicity) &
An under-declared footprint, a conformance failure of that implementation;
the holdout challenge finds one for the encoder and the stretcher; the strict
profile gives every measured declaration whole-asset dependency \\
\addlinespace
Strips the manifest and redistributes &
State becomes \textsc{unverified} &
Detection of the stripping as intentional; re-signing under a fresh, dishonest
chain \\
\addlinespace
Modifies the asset after signing &
Hard binding fails; state becomes \textsc{invalid} &
Attacks on the cryptographic layer itself \\
\addlinespace
Signs a false assertion with a valid key &
Nothing. The contract governs propagation, not truthfulness &
Everything: a valid signature proves signing, not truth \\
\addlinespace
Records a loudspeaker playing synthetic speech &
A captured claim propagates, correct at the sensor layer &
The semantic origin; see the replay limitation in the manuscript \\
\bottomrule
\end{longtable}

\section{EM assertion schema}
\label{s:schema}

The assertion label is \texttt{\Clabel{}}. The payload is:

\begingroup\small
\begin{verbatim}
{
  "schema": "https://github.com/1D-Coda/em-audio/blob/main/docs/em-audio-schema-1.0.md",
  "policy": "complete-source",
  "operator": "trim",
  "operatorParameters": { "start": 4800, "end": 43200 },
  "asset": { "sampleRate": 16000, "sampleCount": 38400 },
  "intervals": [
    {
      "regionOfInterest": {
        "type": "temporal",
        "time": { "type": "npt", "start": "0.000000000", "end": "0.612500000" }
      },
      "samples": { "start": 0, "end": 9800 },
      "provenance": ["C"],
      "state": "CAPTURED",
      "support": { "capture-support": 0.9 },
      "applicability": { "capture-support": ["digital-asset"] },
      "lineage": ["urn:emaudio:librispeech:1272-128104-0000"]
    },
    {
      "regionOfInterest": {
        "type": "temporal",
        "time": { "type": "npt", "start": "0.612500000", "end": "1.437500000" }
      },
      "samples": { "start": 9800, "end": 23000 },
      "provenance": ["G"],
      "state": "GENERATED",
      "support": { "capture-support": 0.1 },
      "applicability": { "capture-support": ["digital-asset"] },
      "lineage": ["urn:emaudio:espeak-ng:1.52.0:0042"]
    },
    {
      "regionOfInterest": {
        "type": "temporal",
        "time": { "type": "npt", "start": "1.437500000", "end": "2.400000000" }
      },
      "samples": { "start": 23000, "end": 38400 },
      "provenance": ["C"],
      "state": "CAPTURED",
      "support": { "capture-support": 0.9 },
      "applicability": { "capture-support": ["digital-asset"] },
      "lineage": ["urn:emaudio:librispeech:1272-128104-0000"]
    }
  ]
}
\end{verbatim}
\endgroup

Normative fields are \texttt{provenance} (\texttt{null} means unverified and is
never an atom), \texttt{support}, \texttt{applicability} and \texttt{lineage}.
The \texttt{state} field is a presentation label. A channel omitted from
\texttt{support} is unavailable, which is different from a channel present with
a low value. Temporal ranges follow the C2PA \texttt{npt} convention with
inclusive start and exclusive end; sample-exact bounds are carried alongside
because \texttt{npt} strings are lossy at sample resolution.

\paragraph{Dependency declaration} Schema~1.0 assertions may carry an optional
\texttt{dependencyDeclaration} object: \texttt{version}, \texttt{profile}
(\texttt{declared} or \texttt{strict}), \texttt{build} (the processing build the
footprints were calibrated on), \texttt{outputSampleCount}, \texttt{sources}
(each with its \texttt{sampleCount}) and \texttt{pieces}, each with its
\texttt{output} and \texttt{sourceRange} spans, \texttt{footprint},
\texttt{basis} (\texttt{analytical}, \texttt{measured} or
\texttt{whole-asset}) and \texttt{label}, and \texttt{footprintTable}, the
module path and SHA-256 digest of the footprint table the declaration was
computed from. A consumer rebuilds the map from it and
recomputes the intervals (\path{em_audio/consumer.py}); assertions without it
remain valid and are reported as unverifiable.
For the trim above, under the strict profile, the declaration is as follows; the
digest and the recorded FFmpeg version header are the exact values of the
evaluated release.

\begingroup\footnotesize
\begin{verbatim}
"dependencyDeclaration": {
  "version": "1.0", "profile": "strict",
  "build": "ffmpeg version 9.0.1 Copyright (c) 2000-2026 the FFmpeg developers",
  "footprintTable": {
    "module": "em_audio/operators.py",
    "sha256":
      "f712c10710234527e1a10a06e40943cdce2bda64dae205259b1faa33f26c2f40" },
  "outputSampleCount": 38400,
  "sources": { "clip": { "sampleCount": 48000 } },
  "pieces": [ { "output": { "start": 0, "end": 38400 }, "source": "clip",
                "sourceRange": { "start": 4800, "end": 43200 },
                "footprint": 0, "basis": "analytical", "label": "trim" } ]
}
\end{verbatim}
\endgroup

The \texttt{build} field records only the first line of \texttt{ffmpeg
-version}. That header does not distinguish builds that share a version number,
which the independent runs show can differ, and the verifier does not compare
it: exact build identity is not enforced in this release. A deployment that
needs it should record the executable's digest and the output of
\texttt{ffmpeg -buildconf} alongside the declaration.

\texttt{"clip"} is a local name for a source within the declaration. The
consumer binds it to the source's own evidence through the derived asset's
\texttt{parentOf} ingredient, whose stored manifest carries the source's
interval assertion. A composition would bind each name to one
\texttt{componentOf} ingredient in the same way; the signed workflow evaluated
here covers the single-source case only.

The \texttt{schema} value is an identifier, not an immutable reference: it is the
string stored in every signed fixture and is left unchanged so those signatures
remain valid. The authoritative schema for the evaluated results is the copy of
\path{docs/em-audio-schema-1.0.md} archived with release \Rtag{} (the Zenodo
record cited in the manuscript's Code availability statement).

\section{Composition manifest wiring}
\label{s:composition}

The C2PA-native composition fixtures of manuscript Section~\SecComposition{}
record each source as a \texttt{componentOf} ingredient carrying its full
manifest, with a \texttt{c2pa.placed} action whose hashed-URI reference the
claim generator resolves at signing time. A frozen composition manifest with
both ingredients, both resolved \texttt{c2pa.placed} references, both source
\texttt{digitalSourceType} values and the EM aggregate is committed at
\path{fixtures/expected/manifest_composition_componentOf.json}. Wiring pitfalls
are listed in \path{docs/C2PA_TROUBLESHOOTING.md}.

\section{Robustness-arm materials}
\label{s:robustness}

The neural arm uses Piper \texttt{piper-tts} 1.7.0 (MIT, VITS architecture)
with the voice \path{en_US-ljspeech-medium}, whose training data is the LJ
Speech Dataset. That dataset states ``This dataset is in the public domain in
the US (and most likely other countries as well)'' and consists of LibriVox
recordings by Linda Johnson of non-fiction texts published between 1884 and
1964, all public domain. The \texttt{lessac} voices were considered and
rejected: their Blizzard 2013 training data carries a restrictive licence. The
arm uses that pretrained voice as distributed; no voice is trained or adapted
here, and no other speaker's recordings are used.

\begingroup
% Line breaking inside a digest is fine, a page break inside one is not, so
% this paragraph is kept on a single page.
\interlinepenalty=10000
The model is fetched rather than committed, by \path{tools/fetch_voice.sh},
which verifies both files against pinned digests. For
\path{en_US-ljspeech-medium.onnx} the SHA-256 is
\digest{6f52a751e2349abe7a76735eb09dc1875298c77ea2342ffd2fef79ff81b87f22}, and
for \path{en_US-ljspeech-medium.onnx.json} it is
\digest{141d612cc0a95ed7efc1ca936b845c2364967f2e9217c5dbfcf69fc4d6c65860}.
Voice output is 22\,050~Hz and is resampled to 16~kHz by stock ffmpeg before
use.
\par\endgroup

The noise arm generates pink noise locally, seeded by clip index so the arm is
deterministic, and mixes it under the clip at $-18$~dB. The exact command, as
executed:

\begin{verbatim}
ffmpeg -i <clip> \
  -f lavfi -i "anoisesrc=r=16000:colour=pink:amplitude=0.12:seed=<index>" \
  -filter_complex "[1:a]volume=-18dB[nz];\
                   [0:a][nz]amix=inputs=2:duration=first:normalize=0[out]" \
  -map "[out]" -c:a pcm_s16le <output>
\end{verbatim}

The noise is modelled as a second generated-derived source covering the whole
clip, so under the contract every captured segment becomes mixed ancestry while
the generated segment stays generated, union adding no new atom. Both facts are
asserted per interval by the experiment.

\section{Reproduction}
\label{s:repro}

\begin{alltt}
git clone https://github.com/1D-Coda/em-audio \bs
  && cd em-audio \bs
  && git checkout \Rtag{} \bs
  && python3 tools/fetch_partialspoof.py \bs
  && python3 tools/fetch_partialedit.py \bs
  && ./run_all.sh
\end{alltt}

\path{run_all.sh} fetches the checksum-pinned corpus, generates the test
signing credential, runs the named regression tests and the experiments, and prints every step it
skipped (the PartialSpoof and PartialEdit runs execute only when those subsets
have been fetched; each fetch keeps about 70~MB); it then regenerates every table, macro and figure from the machine-readable results,
writes \path{results/PREFLIGHT.txt}, and exits non-zero if any conformance
check fails. The preflight report is the single source of every number in the
manuscript; a checker script verifies that no numeric result appears in the
manuscript source outside the generated macro file.

\paragraph{What to record} The environment: the person or group, the machine,
operating system, hardware, runtime and tool versions, and the release tag and
commit reproduced. The execution: a clean clone, dependency installation, a full
\path{run_all.sh}, and a comparison of the emitted machine-readable results
against those shipped with that tag.

\paragraph{Which outputs must match, and which must not} The classification
matters more than the checklist, and it is applied by
\path{tools/verify_reproduction.py} rather than by judgement after the fact.

The \emph{deterministic} outputs are those a correct reimplementation must match
exactly: conformance counts and pass/fail totals, promotion classifications,
lineage sets, interval structure, support-channel behaviour, and the
before-and-after-signing identity of every decoded-essence digest. Any
difference among them exits non-zero. The digests themselves are not compared
across environments: a different decoder build decodes to different bytes, so
equality of a digest is defined within one environment, and what must transfer
is that signing leaves it unchanged.

The \emph{environment-dependent} outputs are those that should be expected to
differ, and whose difference is not a defect: wall-clock timings and their
interquartile ranges, the absolute cost figures of Section~\SecOverhead{} of the
manuscript, and the bytes of a freshly signed file, which differ because signing
incorporates a timestamp. A reproduction reporting identical timings would be
the surprising result, not a reassuring one.

A difference in the first class should be classified rather than absorbed, as an
environment difference, a tool-version difference, an ambiguity in the
specification of the contract, or a genuine discrepancy, and reported as such
whichever it turns out to be.

\section{Independent reproduction: environment and per-operator detail}
\label{s:independent}

The run reported in Section~\SecIndependent{} of the manuscript was performed by
D.~A.~Balderrama-Alvarez (Universidad de Sonora, ORCID 0009-0002-5180-0406),
who is not an author. His environment, as his own preflight report records it:
a \IRcpu{} with \IRcores{} logical processors and \IRmem{}~GiB, \IRos{} on
\IRarch{}, Python~\IRpython{}, FFmpeg~\IRffmpeg{},
\texttt{c2patool}~\IRcTwopatool{}, Node~\IRnode{} and
eSpeak~NG~\IRespeak{}. The reference machine for every other
number in the paper is \Vos{} on \Varch{} with Python~\Vpython{},
FFmpeg~\Vffmpeg{}, \texttt{c2patool}~\VcTwopatool{} and Node~\Vnode{}. He ran the
package twice on consecutive days. The first run was incomplete: two Python
dependencies were missing from the instructions he was given, so the figure and
robustness arms did not execute. The second run, on \IRdate{}, is the one
released and tabulated here; both runs returned the same values for the two
operators whose declarations his build does not satisfy. His unmodified outputs, his preflight report, his full pipeline log and
the classified comparison his run produced are released in
\path{results/independent/}.

Two later runs on earlier released packages are summarised in Section~\SecIndependent{} of the
manuscript; their environments are as follows. B.~C.~Guerra Flores ran it on
an \XBMacCpu{} under \XBMacOs{}, with FFmpeg~\XBMacFfmpeg{} from the same
package manager as the reference machine on different hardware, and with a
\texttt{c2patool} and a Node newer than the reference. Her
\texttt{verify\_reproduction.py} exited zero and \XBMacVerdict{} matched the
package she was given, an earlier release whose suite had \XBMacChecks{}
checks, the \XBMacEnvDiffs{} differences being timings.
D.~A.~Balderrama-Alvarez ran the released package, tag~\XBWinTag{}, as a
single unattended pass on \XBWinOs{} with FFmpeg~\XBWinFfmpeg{}, a different
build of the same FFmpeg version. His verifier reported \XBWinDetDiffs{}
deterministic differences, all of them in the MP3 kernel measurements of
the containment probes and in the time-stretch model deviation, on which the two builds
also disagree, \XBRefStretchDev{} against \XBWinStretchDev{} samples, both
inside the \XBStretchGuard{}-sample margin.

\paragraph{Run history} His first run was incomplete because two dependencies
were missing from the package instructions, since fixed; the second, reported in
the manuscript, gave the same values for both findings. The package he ran,
tag~\IRtag{}, predates the evaluated release \Rtag{}; neither finding's value
changed between them, which we checked. \path{tools/verify_reproduction.py}
returned \texttt{REPRODUCTION INCOMPLETE} with \IRdet{} differences, and
\texttt{run\_all.sh} ended in \texttt{RUN FAILED} on two packaging defects that
touch no reported number. His suite had \IRachecks{} checks, all passing; the
released suite has \Achecks{} because structural checks were added after his
run, and the Windows run of tag~\XBWinTag{} reproduces the \XBWinChecks{} exactly.
On his build \texttt{aselect} retains \IRsilLen{} samples where the reference
retains \IRsilLenRef{}. The Windows build's time-stretch model deviation is
\XBWinStretchDev{} samples against \XBRefStretchDev{} on the reference, both inside
the \XBStretchGuard{}-sample margin, and the declaration of \IRmpDeclared{}
samples contained every macOS and Windows measurement (\XBMacOutside{} and
\XBWinOutside{} probes outside).

\paragraph{Cost on his machine} His bookkeeping cost is \IRabs{}~ms per audio
minute against \IRabsRef{} on the reference machine, a difference of
\IRabsGap{}\%; the ratio to the baseline is \IRratio{} against \IRratioRef{}, a
difference of \IRratioGap{}\%. Across \MstabRepeats{} processes in his sitting the
absolute figure varied by \IRabsCv{}\% and the ratio by \IRratioCv{}\%. One pair of
machines shows the ratio surviving one change of hardware; it does not bound how
far it can move.

Table~\ref{tab:independent} gives the containment result per operator on both
machines rather than only the operator that differs, so the selection made in
the main text can be checked. \IRopsAgree{} of the \IRopsNonzero{} operators
that declare a non-zero footprint reproduce their measured reach exactly across
the two builds; the MP3 encoder does not, and its declared footprint does not
contain the reach his build measures. The other \IRopsZero{} rows declare no
footprint and measure no reach on either machine.

\begin{table}[htbp]
\centering
\caption{Kernel-support containment on the reference machine and on the
independent reproducer's machine. \emph{Declared} is the footprint of
Table~\TabOperators{} of the manuscript, in source samples; \emph{Reach} is the widest
measured influence of a single input sample; \emph{Outside} counts output
samples whose dependency falls outside the declared support. The dagger marks
the one operator whose reach did not reproduce.}
\label{tab:independent}
\small\setlength{\tabcolsep}{4.5pt}
\begin{tabular}{lrrrrr}
\toprule
Operator & Declared & Reach (ref.) & Reach (ind.) & Outside (ref.) & Outside (ind.) \\
\midrule
retained-run selection & 4\,096 & 3\,277 & 3\,277 & 0 & 0 \\
time stretch 1.10 & 2\,577 & 1\,906 & 1\,906 & 0 & 0 \\
transcode to MP3$^{\dagger}$ & 2\,304 & 1\,555 & 4\,317 & 0 & 110 \\
resample 16 to 8 kHz & 97 & 32 & 32 & 0 & 0 \\
amplitude normalisation & 0 & 0 & 0 & 0 & 0 \\
transcode to FLAC & 0 & 0 & 0 & 0 & 0 \\
trim to the middle 80\% & 0 & 0 & 0 & 0 & 0 \\
\bottomrule
\end{tabular}

\end{table}

\section{Holdout challenge: per-operator detail}
\label{sec:supp-holdout}

The holdout challenge probes each operator \KtwoProbesPerOp{} times: \KtwoContexts{}
signal contexts, \KtwoHoldoutContexts{} of them (white noise and a speech-like
voice) unused while the declarations were set, \KtwoAmps{} impulse amplitudes
from \KtwoAmpMin{} to \KtwoAmpMax{}, and \KtwoPositions{} probe positions
disjoint from those of the containment probes, all fixed before the run. Its outputs are in
\path{results/machine_readable/K2_footprint_holdout.json}, one row per probe.

The MP3 encoder exceeded its declaration in \KtwoMpExceed{} probes, with a reach
of \KtwoMpReach{} source samples, and the time stretcher in \KtwoStExceed{},
with a reach of \KtwoStReach{}; the stretcher's exceedances occur at every
amplitude down to \KtwoStMinAmp{}. Resampling (reach \KtwoResampleReach{}),
retained-run selection (reach \KtwoSelReach{}) and fixed-gain normalisation
showed none. With the gain estimated from each input, normalisation reached
\KtwoNormCtrlReach{} source samples in \KtwoNormCtrlExceed{} probes: the whole
input, through the gain, which the contract treats as a control dependency
rather than represented content. A finite challenge can refute a declaration
and cannot establish one, so the absence of an exceedance here does not show
containment at every input.
Under the strict profile, which gives every measured declaration whole-asset
dependency and fits each emitted map to the decoded output length,
\KtwoStrictProbes{} probes of any content operator fell outside. The promotion
demonstrations, one per failing operator, are stored under
\texttt{promotion\_demonstration} in the same file.

\section{Raw rows behind two summary figures}

Tables~\ref{tab:benchrows} and~\ref{tab:calrows} report the process-level timing measurements and the
per-operator calibration summaries. Both are emitted by
\path{tools/make_tables.py} from the same result files the main text reads.

\begin{table}[htbp]
\centering
\caption{The benchmark of Section~\SecOverhead{} as nine independent processes. The
coefficient of variation is the sample standard deviation over the mean, in
percent, with Bessel's correction. The absolute columns move with machine state;
the last row gives each column's coefficient of variation.}
\label{tab:benchrows}
\begin{tabular}{lrrr}
\toprule
Process & EM (ms/audio-min) & Baseline (ms/audio-min) & Ratio \\
\midrule
1 & 1.4179 & 0.4083 & 3.4725 \\
2 & 1.3753 & 0.4128 & 3.3315 \\
3 & 1.3644 & 0.4106 & 3.3233 \\
4 & 1.3636 & 0.4125 & 3.3059 \\
5 & 1.4081 & 0.4115 & 3.4216 \\
6 & 1.4505 & 0.4150 & 3.4954 \\
7 & 1.3688 & 0.4101 & 3.3379 \\
8 & 1.3738 & 0.4089 & 3.3599 \\
9 & 1.3989 & 0.4186 & 3.3422 \\
\midrule
CV (\%) & 2.143 & 0.780 & 2.046 \\
\midrule
\multicolumn{4}{l}{\footnotesize measured 2026-10-06T10:30:51Z UTC, load average [9.55, 10.42, 11.65]} \\
\bottomrule
\end{tabular}

\end{table}

\begin{table}[htbp]
\centering
\caption{Calibration probes per operator, with the context and source position
that produced the widest measured influence. The full per-probe record, all
\Kprobes{} rows, is in \texttt{results/machine\_readable/CALIBRATION.json};
this table gives the worst case for each operator,
which is the value the declaration has to cover. The \Kops{} rows are the
single-source configurations: concatenation assembles the corpus rather than
transforming it, and overlay takes several sources at once.}
\label{tab:calrows}
\begin{tabular}{lrrlrr}
\toprule
Operator & Probes & Max reach & Worst context & At sample & Declared \\
\midrule
retained-run selection & 40 & 3\,277 & tone & 6\,569 & 4\,096 \\
time stretch 1.10 & 40 & 1\,906 & tone & 6\,537 & 2\,577 \\
transcode to MP3 & 40 & 1\,555 & dense & 8\,192 & 2\,304 \\
resample 16 to 8 kHz & 40 & 32 & tone & 1\,638 & 97 \\
trim to the middle 80\% & 40 & 0 & tone & 64 & 0 \\
transcode to FLAC & 40 & 0 & tone & 64 & 0 \\
amplitude normalisation & 40 & 0 & tone & 64 & 0 \\
\bottomrule
\end{tabular}

\end{table}

\end{document}
